\documentclass[final]{beamer}

% Linguistic MDL Explanations with LLMs — via Occam-Style Prompt Discovery
% 24in x 36in portrait, 3 columns
% Compile with: lualatex (via latexmk, engine set in .latexmkrc)

\usepackage[size=custom,width=60.96,height=91.44,scale=0.95]{beamerposter}
\usetheme{gemini}

\usepackage{graphicx}
\usepackage{booktabs}
\usepackage{amsmath, amssymb, bm}
\usepackage{multicol}
\usepackage{array}
\usepackage{ragged2e}
\usepackage{caption}
\usepackage{xcolor}
\usepackage{tikz}
\usetikzlibrary{positioning}
% enumitem is incompatible with beamer; use beamer's native list spacing instead

% -----------------------------
% Gemini-like color palette
% -----------------------------
\definecolor{GeminiBlue}{RGB}{33, 97, 140}
\definecolor{GeminiLightBlue}{RGB}{236, 244, 252}
\definecolor{GeminiDark}{RGB}{36, 41, 46}
\definecolor{GeminiGray}{RGB}{90, 98, 108}
\definecolor{GeminiAccent}{RGB}{0, 114, 178}

\setbeamercolor{headline}{fg=white,bg=GeminiBlue}
\setbeamercolor{title in headline}{fg=white}
\setbeamercolor{author in headline}{fg=white}
\setbeamercolor{institute in headline}{fg=white}

\setbeamercolor{block title}{fg=GeminiBlue,bg=white}
\setbeamercolor{block body}{fg=black,bg=white}

\setbeamercolor{block alerted title}{fg=white,bg=GeminiBlue}
\setbeamercolor{block alerted body}{fg=black,bg=GeminiLightBlue}

\setbeamercolor{block example title}{fg=white,bg=GeminiAccent}
\setbeamercolor{block example body}{fg=black,bg=GeminiLightBlue}

\setbeamercolor{normal text}{fg=GeminiDark,bg=white}
\setbeamercolor{footline}{fg=white,bg=GeminiBlue}

% -----------------------------
% Layout
% -----------------------------
\newlength{\sepwidth}
\newlength{\colwidth}
\setlength{\sepwidth}{0.025\paperwidth}
\setlength{\colwidth}{0.30\paperwidth}

\setlength{\parskip}{0.4em}
\renewcommand{\arraystretch}{1.2}

\captionsetup{font=large}
\setlength{\leftmargini}{1.5em}
\setlength{\leftmarginii}{1.5em}

% -----------------------------
% Metadata
% -----------------------------
\title{Linguistic MDL Explanations with LLMs \\ \large via Occam-Style Prompt Discovery}
\author{Yassine Meliani$^1$ \and George Kontorousis$^1$ \and Eric Elmoznino$^2$}
\institute{$^1$Department of Computer Science, McGill University \quad\textbar\quad $^2$Mila Québec}

% -----------------------------
% Helper commands
% -----------------------------
\newcommand{\panelspace}{\vspace{0.6em}}
\newcommand{\placeholderfigure}[2][0.95\linewidth]{%
  \begin{center}
    \fbox{\parbox[c][#2][c]{#1}{\centering Placeholder figure}}
  \end{center}
}

\begin{document}

\begin{frame}[t]
\begin{columns}[t]

\begin{column}{\sepwidth}\end{column}

% =============================
% Column 1
% =============================
\begin{column}{\colwidth}

\begin{alertblock}{Takeaway}
\justifying
We cast linguistic explanation discovery as a \textbf{two-part MDL objective}: good explanations must be both \textbf{short} and \textbf{predictive}. This gives a principled alternative to heuristic prompt search and exposes trivial explanations that fit the data without explaining it.
\end{alertblock}

\begin{block}{Why This Problem?}
\justifying
Prompted LLMs can act as natural-language programs [1], but existing explanation-discovery methods are mostly heuristic.

\begin{itemize}
  \item Manual prompting and unconstrained generation do not define what makes an explanation \emph{good}
  \item Long descriptions can memorize the dataset; short descriptions can miss predictive structure
  \item We use Occam's razor directly: prefer explanations that are concise \emph{and} predict the observed data
\end{itemize}
\end{block}

\begin{block}{MDL Objective}
\justifying
An explanation $z$ is a natural-language hypothesis for dataset $x$. A frozen evaluator $p^*$ scores $z$ with a two-part code:

\begin{equation}
\mathrm{MDL}(z, x) = \underbrace{-\log p^*(z)}_{\text{explanation complexity}} + \underbrace{-\log p^*(x \mid z)}_{\text{data complexity given } z}
\end{equation}

\begin{itemize}
  \item $-\log p^*(z)$ penalizes verbose or implausible explanations
  \item $-\log p^*(x \mid z)$ rewards explanations that compress the data
  \item Minimizing MDL is equivalent to maximizing the posterior $p^*(z \mid x)$
  \item A token mask $m$ restricts scoring to prediction-relevant output spans
\end{itemize}
\end{block}

\begin{block}{What We Built}
\justifying
\begin{itemize}
  \item \textbf{Policy} $p_\theta(z \mid x)$: generates candidate natural-language explanations
  \item \textbf{Evaluator} $p^*$: frozen LLM scoring prior NLL $-\!\log p^*(z)$ and likelihood NLL $-\!\log p^*(x \mid z)$
  \item \textbf{Training signal}: RL reward derived from the MDL objective
  \item \textbf{Amortization}: after training, one forward pass proposes an explanation for a new dataset
\end{itemize}
\end{block}

\end{column}

\begin{column}{\sepwidth}\end{column}

% =============================
% Column 2
% =============================
\begin{column}{\colwidth}

\begin{block}{Training Pipeline}
\justifying
\begin{center}
\large
Dataset $x$ $\rightarrow$ Policy $p_\theta(z \mid x)$ $\rightarrow$ Evaluator $p^*$ $\rightarrow$ Reward $R$ $\xrightarrow{\mathrm{RL}}$ update $p_\theta$
\normalsize
\end{center}

\begin{itemize}
  \item The policy proposes an explanation $z$
  \item The evaluator scores both its \textbf{complexity} and its \textbf{predictive fit}
  \item The reward updates only $p_\theta$; the evaluator stays frozen
  \item This discourages reward hacking and keeps the objective tied to compression
\end{itemize}
\end{block}

\begin{block}{Reward Signal}
\justifying
The policy is trained with an MDL-derived reward:
\begin{equation}
R(z, x, m;\, p^*) = -\log p^*(z) - \log p^*(x \mid z)
\end{equation}
\begin{itemize}
  \item REINFORCE/PPO optimizes $p_\theta$ under this reward
  \item The frozen evaluator defines a stable objective throughout training
  \item Low MDL explanations are short, reusable, and predictive
\end{itemize}
\end{block}

\begin{alertblock}{Main Evidence: Bayesian Scoring Rejects Trivial Explanations}
\justifying
Sequence: $[1, 1, 2, 5, 12, 27, 58, 121, 248, 503]$

\begin{center}
\scriptsize
\resizebox{\linewidth}{!}{%
\begin{tabular}{lcccc}
\toprule
Explanation & Cond.\ $\downarrow$ & MDL $\downarrow$ & Prior & Lik. \\
\midrule
Generated: $a_n = a_{n-1} + 2a_{n-2}$ & 2.37 & 67.66 & 26.87 & 40.79 \\
Concise: $2^n - n$                      & 39.35 & 81.97 & 64.26 & 17.71 \\
Dummy: repeat sequence                 & 63.88 & 106.50 & 106.50 & 0.00 \\
False: $3n + 1$                         & 42.27 & 120.29 & 61.05 & 59.24 \\
\bottomrule
\end{tabular}
}
\end{center}

\begin{itemize}
  \item The dummy explanation achieves perfect likelihood by restating the sequence, but incurs a very high prior penalty
  \item The concise closed-form rule is short but predicts the sequence less well than the correct recurrence
  \item The correct recurrence attains the best overall two-part MDL score
  \item MDL scoring distinguishes \emph{explaining} the data from merely \emph{describing} it
\end{itemize}
\end{alertblock}

\begin{block}{Benchmarks}
\justifying
\textbf{We evaluate the framework on two task families:}
\begin{itemize}
  \item \textbf{WorldLLM} [2]: procedural text environments where the target is the next-state transition
  \item \textbf{AutoDiscovery} [3]: real-world tabular datasets where the target is the output variable
  \item In both settings, masking restricts the likelihood term to the prediction target rather than the full prompt
\end{itemize}

\textbf{Input / target format}
\begin{itemize}
  \item WorldLLM: $\langle$state, action$\rangle \rightarrow \texttt{[NEXT-STATE TARGET]}$
  \item AutoDiscovery: $(v_1, v_2, \ldots) \rightarrow \texttt{[OUTPUT TARGET]}$
\end{itemize}
\end{block}

\end{column}

\begin{column}{\sepwidth}\end{column}

% =============================
% Column 3
% =============================
\begin{column}{\colwidth}

\begin{block}{Baselines}
\begin{itemize}
  \item \textbf{MDL Policy (ours)}: amortized policy trained to generate explanations under the full two-part MDL code
  \item \textbf{Standard ICL}: no explicit explanation; predictions come directly from in-context examples
  \item \textbf{Monte Carlo Search}: per-dataset candidate search under the same MDL objective but without a learned policy
\end{itemize}
\end{block}

\begin{block}{Evaluation Criteria}
\justifying
\begin{itemize}
  \item \textbf{Training compression}: two-part code length on the observed dataset
  \item \textbf{Held-out predictive loss}: how well an explanation transfers to unseen examples
  \item \textbf{Explanation length}: a simple interpretability proxy
  \item \textbf{Search cost}: amortized policy generation versus per-dataset search
\end{itemize}
\end{block}

\begin{block}{What This Poster Demonstrates}
\begin{itemize}
  \item A principled \textbf{MDL formulation} for linguistic explanation discovery
  \item An \textbf{RL training pipeline} that learns explanations rather than searching from scratch for each dataset
  \item A concrete cas study showing that \textbf{Bayesian two-part scoring} penalizes trivial explanations that likelihood-only scoring prefers
\end{itemize}
\end{block}

\begin{block}{Conclusion}
\begin{itemize}
  \item MDL gives a clean objective for discovering linguistic explanations with LLMs
  \item The two-part score captures a useful trade-off between explanation simplicity and predictive adequacy
  \item Prior-aware scoring is preferable to likelihood-only scoring when trivial explanations are available
  \item The framework supports amortized explanation discovery across interactive and scientific datasets
\end{itemize}
\end{block}

\begin{block}{References}
\tiny
[1] Petrov et al. Prompting as a Universal Approximator. arXiv, 2024.\\
[2] Levy et al. WorldLLM. arXiv, 2025.\\
[3] Agarwal et al. AutoDiscovery. arXiv, 2026.\\
[4] Ellis et al. Human-like Few-Shot Learning via Bayesian Reasoning over Natural Language. NeurIPS, 2023.\\
[5] Hu et al. Amortizing Intractable Inference in Large Language Models. ICLR, 2024.
\end{block}

\end{column}

\begin{column}{\sepwidth}\end{column}

\end{columns}
\end{frame}

\end{document}
